\documentclass{article}
\usepackage{amsmath,amssymb,hyperref}
\title{Connected comma categories and selective right-Kan parallel}
\date{15 November 2026}
\begin{document}\maketitle
\paragraph{Scope.} Independently checked categorical criterion for separated binary assembly of selected comparison adjoints. It is not parallel on all lax-action arrows, nor a theorem about arbitrary correlated/coupled inputs. Pointwise Kan extensions and limits are established machinery; no novelty assertion is made.

\paragraph{Connection to shared-index restriction.} For a functor $i:B\to A$, the conjoint in the covariant-target module convention has components $i^*(b,a)=A(a,ib)$. Its singleton action on covariant presheaves is restriction $i^*:[A,\mathrm{Set}]\to[B,\mathrm{Set}]$, by co-Yoneda. The right adjoint supplied by the action enlargement is therefore the ordinary right Kan extension $\operatorname{Ran}_i$. Thus an inclusion of shared into independent enriched cone presentations has a concrete meaning in the ordinary module realization, rather than an unspecified new adjoint.

\paragraph{Canonical separated comparison.} Given $i:B\to A$, $j:D\to C$, $P:B\to\mathrm{Set}$ and $Q:D\to\mathrm{Set}$, there is a canonical map
\[
 (\operatorname{Ran}_i P)\boxtimes(\operatorname{Ran}_j Q)
 \longrightarrow\operatorname{Ran}_{i\times j}(P\boxtimes Q),
\]
where $\boxtimes$ is the external pointwise Cartesian product. At $(a,c)$, set $I=(a\downarrow i)$ and $J=(c\downarrow j)$. The displayed map is
\[
 (\lim_I P)\times(\lim_J Q)\longrightarrow
 \lim_{I\times J}(P\times Q).
\]
Here the diagrams are restricted along the comma projections. The map takes two compatible families to their pair at each pair of indices.

\paragraph{Exact pointwise criterion in Set.} This comparison is invertible for every $P,Q$ precisely when either both $I,J$ are nonempty connected, or both are empty. If both are nonempty connected, compatibility in the $J$ coordinate forces the first component of a joint family to be independent of that coordinate, and compatibility in $I$ does the same for the second component. This constructs the inverse without choices of distinguished objects. If both are empty, both sides are singleton.

For necessity use only constant diagrams. If $I$ is empty and $J$ nonempty, take $Q$ constant two: the source has at least two elements and the target is singleton. The other mixed empty case is dual. If both are nonempty but $I$ has at least two connected components, take $P$ constant one and $Q$ constant two. A target compatible family may vary with the connected component of $I$, while every family in the image is independent of that component. Thus the map is not surjective. Exchange $I,J$ for the remaining case. This proves the criterion even though not every comma diagram need come from $P$ or $Q$; the constant witnesses do come from them.

\paragraph{General positive form.} In any complete category $V$ with Cartesian products, the same nonempty-connected condition is sufficient for the external-product comparison. A constant diagram over a nonempty connected category has limit its constant value, and products preserve limits jointly. These facts identify the two sides by iterated limits. The Set argument supplies sharpness when that realization is required; no sharpness claim for every $V$ is intended.

\paragraph{Stable class and coherence.} Select maps $i$ for which every $a\downarrow i$ is nonempty connected. Equivalently in Set, $\operatorname{Ran}_i$ carries every constant presheaf to that same constant presheaf via its canonical comparison. Identities have the property. It is closed under composition because $\operatorname{Ran}_{ij}\cong\operatorname{Ran}_i\operatorname{Ran}_j$, with the constant comparisons composing, and the criterion detects the property. It is closed under products because the comma categories are products of nonempty connected categories. Iterated separated comparisons are coherent: all send a tuple of compatible families to its componentwise tuple, so the associativity and unit diagrams agree after each limiting projection. The right-Kan composition identifications have the usual coherence from uniqueness of adjoints to restriction. This does not identify right-Kan parallel with arbitrary old-module parallel on mixed arrows.

\paragraph{Nontriviality beyond represented module adjoints.} If $i$ has a left adjoint $L$, every comma $a\downarrow i$ has initial object $(La,\eta_a)$, recovering the earlier opcartesian/coproduct geometry. The condition is weaker. For a nontrivial group $G$, the unique $BG\to1$ has connected nonempty comma category but no initial object. Its right Kan extension is the fixed-point functor $G\text{-}\mathrm{Set}\to\mathrm{Set}$, which need not preserve colimits. For the regular $G$-set, fixed points are empty whereas its orbit quotient is singleton. Independent fixed points nevertheless commute with the external product for independent group actions, as the theorem asserts.

\paragraph{What remains.} This criterion excludes the disconnected-comma duplication responsible for the simplest scalar obstruction, while allowing genuinely noncocontinuous selected adjoints. It does not guarantee commutation with every completed pure module or every left Kan extension. It says nothing automatic about right Kan extension of a coupled presheaf not of the form $P\boxtimes Q$, nor about a shared diagonal among the outputs. For a full relative construction one must derive the relevant comma/connectedness or enriched replacement from the original context/content data, preserve the actual admitted pure interface, and prove the needed mixed comparisons. These remain explicit obligations.
\end{document}
